% =====================================================================
%  LiaScript / Markdown CheatSheet — A4, doppelseitig
%  Workshop "OPAL Schule meets LiaScript", 28.05.2026
%  Begleitmaterial zu Phase 3 (Anwenden) — die Kartentitel entsprechen
%  exakt den "Cheatsheet-Abschnitt:"-Labels in 03_Anwenden_Vorlage.md.
%  Compile:  pdflatex cheatsheet.tex   (zweimal für korrekte Boxen)
% =====================================================================
\documentclass[a4paper,landscape,9pt]{extarticle}

% --- Layout ----------------------------------------------------------
\usepackage[landscape,margin=0.9cm,top=1.0cm,bottom=0.7cm]{geometry}
\usepackage{multicol}
\setlength{\columnsep}{8pt}
\setlength{\parindent}{0pt}
\setlength{\parskip}{0pt}

% --- Sprache & Schrift ----------------------------------------------
\usepackage[T1]{fontenc}
\usepackage[utf8]{inputenc}
\usepackage[ngerman]{babel}
\usepackage{libertine}
\usepackage[varqu,varl,scaled=0.92]{inconsolata}

% --- Farben ----------------------------------------------------------
\usepackage{xcolor}
\definecolor{mdblue}{HTML}{1F4E79}
\definecolor{mdblueLt}{HTML}{E8EEF5}
\definecolor{liagreen}{HTML}{2E6B3D}
\definecolor{liagreenLt}{HTML}{E8F0E9}
\definecolor{codebg}{HTML}{F6F6F4}
\definecolor{rulegray}{HTML}{D0D0D0}
\definecolor{muted}{HTML}{606060}

% --- Pakete ----------------------------------------------------------
\usepackage[most]{tcolorbox}
\usepackage{listings}
\usepackage{qrcode}
\usepackage{fontawesome5}
\usepackage{array}
\usepackage[normalem]{ulem}
\usepackage{hyperref}
\hypersetup{hidelinks}

% --- Listings: Markdown/LiaScript-Quellcode --------------------------
\lstdefinelanguage{lia}{
  morecomment=[l]{\#},
  morecomment=[s]{<!--}{-->},
  sensitive=true
}
\lstset{
  basicstyle=\ttfamily\scriptsize,
  breaklines=true,
  breakatwhitespace=true,
  showstringspaces=false,
  keepspaces=true,
  columns=fullflexible,
  upquote=true,
  literate=%
    {ä}{{\"a}}1 {ö}{{\"o}}1 {ü}{{\"u}}1
    {Ä}{{\"A}}1 {Ö}{{\"O}}1 {Ü}{{\"U}}1
    {ß}{{\ss}}1 {„}{{\glqq}}1 {“}{{\grqq}}1,
  frame=none,
  xleftmargin=0pt,
  aboveskip=0pt,
  belowskip=0pt
}

% --- Karten-Vorlage --------------------------------------------------
%  Zwei Varianten: mdcard (blau, Markdown) und liacard (grün, LiaScript)
%  Innen: links Quellcode, rechts Wirkung — getrennt durch \tcblower.
\tcbset{
  cheatbase/.style={
    enhanced,
    sidebyside,
    sidebyside align=top,
    lefthand width=0.46\linewidth,
    segmentation style={solid, rulegray, line width=0.3pt},
    colback=white,
    boxrule=0.3pt,
    arc=2pt,
    left=4pt,right=4pt,top=2pt,bottom=2pt,middle=4pt,
    before skip=3pt, after skip=3pt,
    breakable=false,
    fonttitle=\sffamily\bfseries\footnotesize,
    coltitle=white
  }
}

\newtcolorbox{mdcard}[1]{cheatbase, colframe=mdblue!75,  colbacktitle=mdblue,   title=#1}
\newtcolorbox{liacard}[1]{cheatbase, colframe=liagreen!75, colbacktitle=liagreen, title=#1}

% Karte für die ganze Spaltenbreite (z. B. Mini-Kurs am Ende)
\newtcolorbox{liawide}[1]{
  enhanced,
  colback=white,
  colframe=liagreen!75,
  boxrule=0.3pt, arc=2pt,
  left=4pt,right=4pt,top=2pt,bottom=2pt,
  before skip=3pt, after skip=0pt,
  breakable=false,
  fonttitle=\sffamily\bfseries\footnotesize,
  coltitle=white, colbacktitle=liagreen, title=#1
}

% Info-Box ohne sidebyside (für Tools/Templates-Spalten unten)
\tcbset{
  infobase/.style={
    enhanced, colback=white,
    boxrule=0.3pt, arc=2pt,
    left=5pt,right=5pt,top=3pt,bottom=3pt,
    before skip=3pt, after skip=0pt,
    breakable=false,
    fonttitle=\sffamily\bfseries\footnotesize, coltitle=white
  }
}
\newtcolorbox{mdinfo}[1]{infobase, colframe=mdblue!75,    colbacktitle=mdblue,   title=#1}
\newtcolorbox{liainfo}[1]{infobase, colframe=liagreen!75, colbacktitle=liagreen, title=#1}

% Kleine Hilfsmacros für die "Wirkung"-Spalte
\newcommand{\prevtxt}[1]{{\small #1}}
\newcommand{\prevsm}[1]{{\footnotesize\color{muted} #1}}

\pagestyle{empty}
\begin{document}

% =====================================================================
% SEITE 1 — MARKDOWN
% =====================================================================
\noindent\begin{tcolorbox}[
  enhanced, frame hidden, sharp corners,
  colback=mdblue, coltext=white,
  left=12pt, right=8pt, top=4pt, bottom=4pt,
  before skip=0pt, after skip=4pt,
  width=\linewidth,
  borderline west={4pt}{-4pt}{mdblue!55}
]
\begin{minipage}[c]{0.66\linewidth}
  {\sffamily\fontsize{26}{28}\selectfont\bfseries Markdown}\hspace{0.4em}%
  {\sffamily\large\color{white!75} — die Basis}\\[2pt]
  {\sffamily\footnotesize\itshape\color{white!90} Was jeder Markdown-Editor versteht — Grundlage für alle LiaScript-Dokumente.}
\end{minipage}\hfill
\begin{minipage}[c]{0.20\linewidth}\raggedright
  {\sffamily\scriptsize\color{white!85}\faGraduationCap~~Phase 3 · Anwenden}\\[2pt]
  {\sffamily\bfseries\normalsize Seite 1\,/\,2}\\[1pt]
  {\sffamily\scriptsize\color{white!75} Markdown-Grundlagen}
\end{minipage}\hfill
\begin{minipage}[c]{0.10\linewidth}\centering
  {\color{black}\colorbox{white}{\hspace{1pt}\qrcode[height=1.55cm,nolinks]{https://liascript.github.io/LiveEditor/}\hspace{1pt}}}\\[1pt]
  {\sffamily\scriptsize\color{white!90}\faEdit~Live-Editor}
\end{minipage}
\end{tcolorbox}

% --- Tag für Aufgabenbezug (klein, oben rechts in der Titelzeile) ---
\newcommand{\aufg}[1]{\hfill{\sffamily\scriptsize\color{white!85}\faTasks~A#1}}

\begin{multicols}{3}

% =====================================================================
% Block A — Textauszeichnung & Struktur (Aufgaben 2-4)
% =====================================================================

% --- A2a. Textauszeichnung ------------------------------------------
\begin{mdcard}{Textauszeichnung\aufg{2}}
\begin{lstlisting}[language=lia]
**Energie**, *kursiv*,
~~durchgestrichen~~,
`Joule`
\end{lstlisting}
\tcblower
\prevtxt{\textbf{Energie}, \textit{kursiv},
\sout{durchgestrichen}, \texttt{Joule}}
\end{mdcard}

% --- A2b. Absätze & Trennlinie --------------------------------------
\begin{mdcard}{Absätze \& Trennlinie\aufg{2}}
\begin{lstlisting}[language=lia]
Erster Absatz.

Zweiter Absatz nach
einer Leerzeile.

---
\end{lstlisting}
\tcblower
\prevtxt{Erster Absatz.\\[3pt]
Zweiter Absatz nach einer Leerzeile.\\[2pt]
\textcolor{rulegray}{\hrulefill}}
\end{mdcard}

% --- A3. Überschriften ----------------------------------------------
\begin{mdcard}{Überschriften\aufg{3}}
\begin{lstlisting}[language=lia]
# Energie
## Worum geht es?
### Energiearten
#### Detail
\end{lstlisting}
\tcblower
{\large\bfseries Energie}\par\vspace{1pt}
{\normalsize\bfseries Worum geht es?}\par\vspace{1pt}
{\small\bfseries Energiearten}\par\vspace{1pt}
{\footnotesize\bfseries Detail}\par\vspace{2pt}
{\scriptsize\color{muted}Bildet die Navigation am rechten Rand.}
\end{mdcard}

% --- A4. Listen -----------------------------------------------------
\begin{mdcard}{Listen\aufg{4}}
\begin{lstlisting}[language=lia]
- Energie definieren
- E_kin und E_pot
  unterscheiden

  - Sonderfall Feder

- Joule benennen

1. erstens
2. zweitens
\end{lstlisting}
\tcblower
\prevtxt{%
$\bullet$ Energie definieren\\
$\bullet$ E\textsubscript{kin} und E\textsubscript{pot} unterscheiden\\[2pt]
\hspace*{1em}$\circ$ Sonderfall Feder\\[2pt]
$\bullet$ Joule benennen\\[3pt]
1. erstens\\
2. zweitens\\[2pt]
{\scriptsize\color{muted}Untergeordnete Listen um \textbf{2 Zeichen} einrücken.}}
\end{mdcard}

% =====================================================================
% Block B — Tabellen, Formeln & Bilder (Aufgaben 5-8)
% =====================================================================

% --- A5. Tabellen ---------------------------------------------------
\begin{mdcard}{Tabellen\aufg{5}}
\begin{lstlisting}[language=lia]
| Energieart | Formel       |
|:-----------|:------------:|
| kinetisch  | 1/2 m v^2    |
| potentiell | m g h        |
| Spann      | 1/2 D s^2    |
\end{lstlisting}
\tcblower
\prevtxt{%
\renewcommand{\arraystretch}{1.05}%
\begin{tabular}{|l|c|}\hline
\textbf{Energieart} & \textbf{Formel}\\\hline
kinetisch  & $\tfrac{1}{2}mv^2$\\
potentiell & $mgh$\\
Spann      & $\tfrac{1}{2}Ds^2$\\\hline
\end{tabular}\\[2pt]
{\scriptsize\color{muted}\texttt{:--} links, \texttt{:--:} zentriert, \texttt{--:} rechts}}
\end{mdcard}

% --- A6. Mathematische Formeln (LaTeX) ------------------------------
\begin{mdcard}{Mathematische Formeln (LaTeX)\aufg{6}}
\begin{lstlisting}[language=lia]
Inline: $E = m c^2$

Block:
$$E_\text{kin}
  = \tfrac{1}{2}\, m\, v^2$$
\end{lstlisting}
\tcblower
\prevtxt{Inline: $E = m c^2$\\[3pt]
Block:\\
$\displaystyle E_\text{kin} = \tfrac{1}{2}\, m\, v^2$\\[3pt]
{\scriptsize\color{muted}Index \texttt{\_}, Bruch \texttt{\textbackslash tfrac}, Quadrat \texttt{\^{}2}.
Ersetzt zerfallene Unicode-Schnipsel.}}
\end{mdcard}

% --- A8. Bilder -----------------------------------------------------
\begin{mdcard}{Bilder\aufg{8}}
\begin{lstlisting}[language=lia]
![Alt-Text](pfad/bild.png)

![Schwingendes Pendel](
  https://.../Pendulum.svg)
\end{lstlisting}
\tcblower
\prevtxt{\fbox{\rule{0pt}{1.2em}\hspace{1.2em}{\scriptsize\color{muted}Bild}\hspace{1.2em}}\\[2pt]
{\scriptsize\color{muted}\textit{Alt-Text in eckigen Klammern — wichtig für Barrierefreiheit. Quelle/Lizenz in der Bildunterschrift nennen.}}}
\end{mdcard}

% =====================================================================
% Block C — Hinweise & Referenz-Karten (Aufgabe 7)
% =====================================================================

% --- A7. Hinweise / Callouts ----------------------------------------
\begin{mdcard}{Hinweise / Callouts\aufg{7}}
\begin{lstlisting}[language=lia]
> Ein einfaches Zitat.

> [!TIP]      gelb
> [!NOTE]     blau
> [!WARNING]  rot
\end{lstlisting}
\tcblower
\prevtxt{%
\textcolor{rulegray}{\rule{1pt}{1.1em}}\hspace{4pt}\textit{Ein einfaches Zitat.}\\[3pt]
\textcolor{mdblue}{\rule{1pt}{1.0em}}\hspace{3pt}{\scriptsize\textbf{TIP} (gelb): Definitionen, Hinweise}\\[1pt]
\hspace*{8pt}{\scriptsize\textbf{NOTE} (blau): Lese-Empfehlungen}\\[1pt]
\hspace*{8pt}{\scriptsize\textbf{WARNING} (rot): Stolperfallen}\\[2pt]
\hspace*{8pt}{\scriptsize Folgezeilen mit \texttt{>} = Inhalt der Box.}}
\end{mdcard}

% --- R1. Aufgabenliste / Checkboxen ---------------------------------
\begin{mdcard}{Checklisten (Referenz)}
\begin{lstlisting}[language=lia]
- [x] erledigt
- [ ] offen
- [ ] noch offen
\end{lstlisting}
\tcblower
\prevtxt{%
\faCheckSquare~erledigt\\
\faSquare~offen\\
\faSquare~noch offen\\[2pt]
{\scriptsize\color{muted}Ideal für eine Lernfortschritts-Bilanz.}}
\end{mdcard}

% --- R2. HTML & Sonderzeichen ---------------------------------------
\begin{mdcard}{HTML \& Sonderzeichen (Referenz)}
\begin{lstlisting}[language=lia]
<u>unterstrichen</u>
<br>  Zeilenumbruch
\* kein Listenpunkt
\end{lstlisting}
\tcblower
\prevtxt{\underline{unterstrichen}\\
{\scriptsize\color{muted}(Zeilenumbruch)}\\
*~kein Listenpunkt\\[2pt]
{\scriptsize\color{muted}Backslash escapt Markdown-Sonderzeichen.}}
\end{mdcard}

\end{multicols}

% --- Untere Zeile: Werkzeuge & Editoren (volle Breite) -------------
\noindent
\begin{mdinfo}{\faToolbox~~Werkzeuge \& Editoren für Markdown}
\footnotesize
\begin{minipage}[t]{0.24\linewidth}
\textbf{\color{mdblue}Lokale Editoren}\\[1pt]
$\bullet$~VS Code \textit{(+ Markdown All in One)}\\
$\bullet$~Obsidian, Typora, Marktext
\end{minipage}\hfill
\begin{minipage}[t]{0.22\linewidth}
\textbf{\color{mdblue}Online \& kollaborativ}\\[1pt]
$\bullet$~LiaScript LiveEditor\\
$\bullet$~hackmd.io \textit{(kollaborativ)}
\end{minipage}\hfill
\begin{minipage}[t]{0.24\linewidth}
\textbf{\color{mdblue}Konvertieren \& Qualität}\\[1pt]
$\bullet$~\texttt{pandoc} → HTML, PDF, DOCX, EPUB\\
$\bullet$~\texttt{markdownlint}, \texttt{prettier}
\end{minipage}\hfill
\begin{minipage}[t]{0.25\linewidth}
\textbf{\color{mdblue}Versionierung}\\[1pt]
$\bullet$~Git \& GitHub/GitLab rendern Markdown direkt\\
$\bullet$~Tabellen: \texttt{tablesgenerator.com/markdown}
\end{minipage}
\end{mdinfo}

\vfill
{\sffamily\scriptsize\color{muted}
\faInfoCircle~Live ausprobieren: \textbf{liascript.github.io/LiveEditor}\hfill
\faGithub~Quellcode: \textbf{github.com/LiaPlayground/Opal\_Schule\_meets\_LiaScript}}

\newpage

% =====================================================================
% SEITE 2 — LIASCRIPT
% =====================================================================
\noindent\begin{tcolorbox}[
  enhanced, frame hidden, sharp corners,
  colback=liagreen, coltext=white,
  left=12pt, right=8pt, top=4pt, bottom=4pt,
  before skip=0pt, after skip=4pt,
  width=\linewidth,
  borderline west={4pt}{-4pt}{liagreen!55}
]
\begin{minipage}[c]{0.66\linewidth}
  {\sffamily\fontsize{26}{28}\selectfont\bfseries LiaScript}\hspace{0.4em}%
  {\sffamily\large\color{white!75} — die Erweiterungen}\\[2pt]
  {\sffamily\footnotesize\itshape\color{white!90} Alles, was Markdown allein nicht kann: Quiz, Animationen, Simulationen, ausführbarer Code, Sprache.}
\end{minipage}\hfill
\begin{minipage}[c]{0.20\linewidth}\raggedright
  {\sffamily\scriptsize\color{white!85}\faGraduationCap~~Phase 3 · Anwenden}\\[2pt]
  {\sffamily\bfseries\normalsize Seite 2\,/\,2}\\[1pt]
  {\sffamily\scriptsize\color{white!75} Erweiterungen \& Interaktion}
\end{minipage}\hfill
\begin{minipage}[c]{0.10\linewidth}\centering
  {\color{black}\colorbox{white}{\hspace{1pt}\qrcode[height=1.55cm,nolinks]{https://github.com/LiaPlayground/Opal_Schule_meets_LiaScript}\hspace{1pt}}}\\[1pt]
  {\sffamily\scriptsize\color{white!90}\faGithub~Workshop-Repo}
\end{minipage}
\end{tcolorbox}

% --- Tag für Aufgabenbezug (grün, oben rechts im Karten-Titel) -----
\newcommand{\aufglia}[1]{\hfill{\sffamily\scriptsize\color{white!85}\faTasks~A#1}}

\begin{multicols}{3}

% =====================================================================
% Block A — Kursrahmen (Aufgabe 1)
% =====================================================================

% --- A1. YAML-Header -------------------------------------------------
\begin{liacard}{YAML-Header\aufglia{1}}
\begin{lstlisting}[language=lia]
<!--
author:   Ihr Name
language: de
narrator: Deutsch Male
import:   https://github.com/
          LiaTemplates/Pyodide
-->
\end{lstlisting}
\tcblower
\prevsm{Steht ganz oben im Dokument.\\
Definiert Autor:in, Sprache, Stimme
und externe Bausteine (\textit{Imports}).\\[2pt]
{\footnotesize Der \texttt{Pyodide}-Import wird für den
Code-Block in Aufgabe 15 gebraucht.}}
\end{liacard}

% =====================================================================
% Block B — Medien & Einbettung (Aufgaben 9-10)
% =====================================================================

% --- A9. Video ------------------------------------------------------
\begin{liacard}{Video\aufglia{9}}
\begin{lstlisting}[language=lia]
!?[Titel](https://www.youtube
.com/watch?v=VIDEO_ID)

!?[Audio](datei.mp3)
\end{lstlisting}
\tcblower
\prevsm{\faYoutube~YouTube, Vimeo, MP4
und MP3 werden direkt eingebettet.\\[2pt]
\textbf{!?} statt \textbf{!} = abspielbar.\\[2pt]
{\footnotesize Auf die Lizenz des Videos achten.}}
\end{liacard}

% --- A10. iframe-Einbettung (PhET-Simulation) -----------------------
\begin{liacard}{iframe-Einbettung\aufglia{10}}
\begin{lstlisting}[language=lia]
??[PhET-Simulation](
  https://phet.colorado.edu/
  .../energy-skate-park_de.html)
\end{lstlisting}
\tcblower
\prevsm{\faWindowMaximize~\textbf{??} erkennt das Format
automatisch (PhET, PDF, GeoGebra, \dots)
und bettet es ein -- kein Tab-Wechsel.\\[2pt]
{\footnotesize Höhe vorab steuern:\\
\texttt{<!-{}- style="height: 700px;" -{}->}\\
direkt über die \texttt{??}-Zeile setzen.}}
\end{liacard}

% =====================================================================
% Block C — Quizze (Aufgaben 11-13)
% =====================================================================

% --- A11. Einfachauswahl --------------------------------------------
\begin{liacard}{Quizformate / Einfachauswahl\aufglia{11}}
\begin{lstlisting}[language=lia]
Welche Einheit hat Energie?

- [( )] Newton
- [(X)] Joule
- [( )] Watt
\end{lstlisting}
\tcblower
\prevsm{Welche Einheit hat Energie?\\[2pt]
$\circ$~Newton\\
$\bullet$~\textbf{Joule}\\
$\circ$~Watt\\[2pt]
{\footnotesize Genau \emph{eine} richtige Antwort: \texttt{[(X)]}.}}
\end{liacard}

% --- A12. Zahleneingabe ---------------------------------------------
\begin{liacard}{Quizformate / Zahleneingabe\aufglia{12}}
\begin{lstlisting}[language=lia]
E_kin für m=4 kg,
v=27 m/s in Joule?

[[1458]]
\end{lstlisting}
\tcblower
\prevsm{$E_\text{kin}$ für $m=4$, $v=27$?\\[2pt]
\fbox{\strut\hspace{5em}}\\[3pt]
{\footnotesize Akzeptiert \texttt{1458}, \texttt{1458.0}, \texttt{1.458e3}.\\
Toleranz: \texttt{[[1458 $\pm$ 1]]}.}}
\end{liacard}

% --- A13. Mehrfachauswahl -------------------------------------------
\begin{liacard}{Quizformate / Mehrfachauswahl\aufglia{13}}
\begin{lstlisting}[language=lia]
Wovon hängt E_pot ab?

- [[X]] Masse m
- [[X]] Höhe h
- [[ ]] Geschwindigkeit v
- [[X]] Gravitation g
\end{lstlisting}
\tcblower
\prevsm{Wovon hängt $E_\text{pot}$ ab?\\[2pt]
\faCheckSquare~Masse $m$\\
\faCheckSquare~Höhe $h$\\
\faSquare~Geschwindigkeit $v$\\
\faCheckSquare~Gravitation $g$\\[2pt]
{\footnotesize \texttt{[[X]]} = beliebig viele richtig.}}
\end{liacard}

% =====================================================================
% Block D — Hilfen & ausführbarer Code (Aufgaben 14-15)
% =====================================================================

% --- A14. Hinweise & Spoiler ----------------------------------------
\begin{liacard}{Hinweise \& Spoiler\aufglia{14}}
\begin{lstlisting}[language=lia]
[[1458]]
- [[?]] Formel:
        1/2 m v^2
*******************

E_kin = 1458 J

*******************
\end{lstlisting}
\tcblower
\prevsm{\faQuestionCircle~Stufe 1: Hinweis \texttt{[[?]]}
(klappbar).\\[3pt]
Stufe 2: Lösungsblock zwischen zwei
\texttt{*****}-Zeilen --- koppelt sich
automatisch an das Quiz darüber.\\[2pt]
{\footnotesize Zweistufige Hilfe = wirksame Binnendifferenzierung.}}
\end{liacard}

% --- A15. Code-Blöcke (Pyodide) -------------------------------------
\begin{liacard}{Code-Blöcke\aufglia{15}}
\begin{lstlisting}[language=lia]
```python  e_kin.py
m = 4; v = 27
E = 0.5 * m * v**2
print(E)
```
@Pyodide.eval
\end{lstlisting}
\tcblower
\prevsm{\colorbox{codebg}{\parbox{0.85\linewidth}{\ttfamily\scriptsize m = 4; v = 27\\E = 0.5*m*v**2\\print(E)}}\\[2pt]
\faPlay~\textbf{Run}~~\textcolor{liagreen}{\texttt{> 1458.0}}\\[2pt]
{\footnotesize \texttt{@Pyodide.eval} direkt darunter setzt den
Run-Button. Setzt den \texttt{Pyodide}-Import (A1) voraus.}}
\end{liacard}

\end{multicols}

% --- Untere Zeile: LiaScript-Templates (volle Breite) ---------------
\noindent
\begin{liainfo}{\faPuzzlePiece~~LiaScript-Templates — im Header per \texttt{import:} einbinden (Auswahl aus \texttt{github.com/LiaTemplates})}
\footnotesize
\begin{minipage}[t]{0.24\linewidth}
\textbf{\color{liagreen}\faTerminal~~Code ausführen}\\[1pt]
$\bullet$~\texttt{Pyodide}~–~Python\\
$\bullet$~\texttt{Skulpt}~–~Python (leichter)\\
$\bullet$~\texttt{Lua}~–~Lua\\
$\bullet$~\texttt{JSCPP}~–~C++\\
$\bullet$~\texttt{Tau-Prolog}~–~Prolog
\end{minipage}\hfill
\begin{minipage}[t]{0.24\linewidth}
\textbf{\color{liagreen}\faChartBar~~Visualisieren}\\[1pt]
$\bullet$~\texttt{mermaid\_template}~–~Diagramme\\
$\bullet$~\texttt{JSXGraph}~–~Mathematik\\
$\bullet$~\texttt{Vega}~–~Datenvisualisierung\\
$\bullet$~\texttt{Tikz-Jax}~–~Tikz-Grafiken\\
$\bullet$~\texttt{ABCjs}~–~Musiknotation
\end{minipage}\hfill
\begin{minipage}[t]{0.24\linewidth}
\textbf{\color{liagreen}\faFlask~~Simulation \& Daten}\\[1pt]
$\bullet$~\texttt{PhET}~–~Physik-Simulationen\\
$\bullet$~\texttt{Wikimedia}~–~Commons-Einbettung\\
$\bullet$~\texttt{AlaSQL}~/~\texttt{SQLite}~–~SQL\\
$\bullet$~\texttt{Random}~–~Quiz-Banken\\
$\bullet$~\texttt{citations}~–~BibTeX/CSL
\end{minipage}\hfill
\begin{minipage}[t]{0.24\linewidth}
\textbf{\color{liagreen}\faGithub~~Übersicht}\\[1pt]
$\bullet$~\texttt{github.com/LiaTemplates}\\
$\bullet$~\texttt{github.com/LiaTemplates/Index}\\
$\bullet$~\texttt{github.com/LiaPlayground}\\
$\bullet$~Live-Editor:\\
\hspace*{0.6em}\texttt{liascript.github.io/LiveEditor}
\end{minipage}
\end{liainfo}

\vspace{4pt}
{\sffamily\scriptsize\color{muted}
\faInfoCircle~Live-Editor: \textbf{liascript.github.io/LiveEditor}\hfill
\faBook~Doku: \textbf{liascript.github.io}\hfill
\faGithub~\textbf{github.com/LiaScript} \quad
CC-BY-SA 4.0 — Sebastian Zug, André Dietrich · OPAL Schule meets LiaScript 2026}

\end{document}
